\documentclass{beamer}

\usetheme{Madrid}

\usepackage{amsmath,amssymb,amsfonts,bm}
\usepackage{physics}

\title{Why Does One Use $\log p(x)$ as a Score Function?}
\author{Morten Hjorth-Jensen}
\date{May 2026}

\begin{document}

%------------------------------------------------
\begin{frame}
\titlepage
\end{frame}

%------------------------------------------------
\begin{frame}{Motivation}
In statistics, machine learning, information theory, and statistical physics, the quantity
\[
\log p(x)
\]
appears everywhere.

\medskip

Examples:
\begin{itemize}
\item Maximum likelihood estimation
\item Fisher information
\item Cross-entropy loss
\item Entropy
\item Bayesian inference
\item Statistical mechanics
\item Variational inference
\end{itemize}

\medskip

Why is the logarithm of the probability so fundamental?
\end{frame}

%------------------------------------------------
\section{Likelihood and Log-Likelihood}

\begin{frame}{Likelihood Function}
Suppose we observe independent data:
\[
x_1,x_2,\dots,x_N.
\]

Assume probability model:
\[
p(x|\theta).
\]

The probability of the full dataset is:
\[
P(D|\theta)
=
\prod_{i=1}^N p(x_i|\theta).
\]

This is called the likelihood.
\end{frame}

%------------------------------------------------
\begin{frame}{Problem with Products}
Products are mathematically inconvenient:
\begin{itemize}
\item numerically unstable,
\item often extremely small,
\item difficult to differentiate,
\item difficult to interpret geometrically.
\end{itemize}

\pause

Taking the logarithm:
\[
\log P(D|\theta)
=
\sum_{i=1}^N \log p(x_i|\theta).
\]

\pause

The logarithm converts products into sums.
\end{frame}

%------------------------------------------------
\begin{frame}{Logarithm Preserves Maxima}
Since the logarithm is monotonic:
\[
\arg\max_\theta P(D|\theta)
=
\arg\max_\theta \log P(D|\theta).
\]

\pause

Thus:
\begin{itemize}
\item maximizing likelihood
\item maximizing log-likelihood
\end{itemize}

are completely equivalent.
\end{frame}

%------------------------------------------------
\section{Score Function}

\begin{frame}{Definition of the Score Function}
The statistical score is defined as:
\[
s(\theta)
=
\partial_\theta \log p(x|\theta).
\]

\pause

Using:
\[
\partial_\theta \log p
=
\frac{1}{p}\partial_\theta p,
\]

the score measures the \emph{relative sensitivity} of the probability distribution.
\end{frame}

%------------------------------------------------
\begin{frame}{Relative Versus Absolute Changes}
Suppose:
\[
p(x) \ll 1.
\]

A small absolute change in probability may still correspond to a large relative change.

\pause

The logarithmic derivative measures:
\[
\frac{\delta p}{p}.
\]

\pause

This is the natural geometry of probability distributions.
\end{frame}

%------------------------------------------------
\section{Fisher Information}

\begin{frame}{Fisher Information}
The Fisher information is:
\[
\mathcal I(\theta)
=
\mathbb E
\left[
(\partial_\theta \log p)^2
\right].
\]

\pause

Interpretation:
\begin{itemize}
\item Measures distinguishability of nearby probability distributions
\item Measures curvature of the likelihood
\item Defines a metric on parameter space
\end{itemize}
\end{frame}

%------------------------------------------------
\begin{frame}{Important Identities}
The score function satisfies:
\[
\mathbb E[\partial_\theta \log p]=0.
\]

\pause

Second derivative identity:
\[
\mathcal I(\theta)
=
-
\mathbb E
\left[
\partial_\theta^2 \log p
\right].
\]

\pause

These elegant identities emerge naturally only when using logarithms.
\end{frame}

%------------------------------------------------
\section{Information Theory}

\begin{frame}{Entropy}
Entropy is defined as:
\[
S
=
-
\sum_x p(x)\log p(x).
\]

\pause

Why does $\log p$ appear?

\medskip

Because logarithms convert multiplicative probabilities into additive information.
\end{frame}

%------------------------------------------------
\begin{frame}{Additivity of Information}
Independent events satisfy:
\[
p(A,B)=p(A)p(B).
\]

\pause

Taking logarithms:
\[
\log p(A,B)
=
\log p(A)+\log p(B).
\]

\pause

Thus information becomes additive.
\end{frame}

%------------------------------------------------
\begin{frame}{Information Content}
Define information:
\[
I(x)=-\log p(x).
\]

\pause

Rare events:
\[
p(x)\ll1
\]

produce large information.

\pause

Common events produce little information.
\end{frame}

%------------------------------------------------
\section{Statistical Mechanics}

\begin{frame}{Boltzmann Distribution}
In statistical mechanics:
\[
p(x)
=
\frac{e^{-\beta E(x)}}{Z}.
\]

\pause

Taking logarithms:
\[
\log p(x)
=
-\beta E(x)-\log Z.
\]

\pause

Therefore:
\[
E(x)
=
-\frac1\beta \log p(x)+\text{const}.
\]

\pause

Energy is proportional to logarithmic probability.
\end{frame}

%------------------------------------------------
\begin{frame}{Connection to Machine Learning}
Many machine learning models are based on logarithmic probabilities:

\begin{itemize}
\item Boltzmann machines
\item Variational autoencoders
\item Diffusion models
\item Bayesian neural networks
\item GAN discriminators
\end{itemize}

\pause

The logarithm naturally appears because probabilities multiply while optimization prefers additive structures.
\end{frame}

%------------------------------------------------
\section{Deep Learning}

\begin{frame}{Cross-Entropy Loss}
Classification networks often use:
\[
L
=
-
\sum_i y_i \log p_i.
\]

\pause

This is the negative log-likelihood of the observed labels.

\pause

Thus cross-entropy follows directly from probability theory.
\end{frame}

%------------------------------------------------
\section{Geometric Interpretation}

\begin{frame}{Geometry of Probability Space}
Probability distributions form a curved manifold.

\pause

The logarithm transforms probabilities into coordinates where:
\begin{itemize}
\item gradients,
\item distances,
\item curvatures,
\item information measures
\end{itemize}

become mathematically tractable.
\end{frame}

%------------------------------------------------
\begin{frame}{Information Geometry}
The Fisher information metric:
\[
ds^2
=
\sum_{ij}
\mathcal I_{ij}
d\theta_i d\theta_j
\]

defines the geometry of statistical models.

\pause

The tangent vectors are:
\[
\partial_\theta \log p(x|\theta).
\]

Thus the score function defines directions in probability space.
\end{frame}

%------------------------------------------------
\section{Summary}

\begin{frame}{Summary}
We use:
\[
\log p(x)
\]

because:

\begin{itemize}
\item probabilities multiply,
\item logarithms convert products into sums,
\item optimization becomes stable,
\item gradients become relative sensitivities,
\item information becomes additive,
\item entropy emerges naturally,
\item statistical mechanics takes exponential form,
\item geometry becomes tractable.
\end{itemize}
\end{frame}

%------------------------------------------------
\begin{frame}{Final Perspective}
The score function:
\[
\partial_\theta \log p(x|\theta)
\]

measures how rapidly the \emph{information content} of an observation changes with model parameters.

\medskip

It therefore lies at the foundation of:
\begin{itemize}
\item statistics,
\item machine learning,
\item Bayesian inference,
\item information theory,
\item statistical physics,
\item modern AI.
\end{itemize}
\end{frame}

%------------------------------------------------
\end{document}
